\section{Methodology: Hidden-Profile Scenario and Questionnaire Instrument}
\subsection{Study Design}

The empirical evaluation of the proactive AI assistant required a group task that resembles the
convergent organizational decisions this thesis addresses and that can generate situations in which
facilitation --- by a human moderator or by an AI agent --- may become relevant.

Houtti et al.\cite{houttiObserveAskIntervene2025} evaluated their virtual co-host with the Lost at Sea
activity, in which participants rank 15 items by importance for survival after a shipwreck.
They selected this task in part because it does not require equal participation and can be completed
individually \cite{houttiObserveAskIntervene2025}.
For the present study, these properties are a limitation.
A ranking of a static, commonly known list does not require participants to pool unique information.
As a result, the discussion is unlikely to produce the information-related difficulties that
characterize many organizational decision meetings, and therefore few moments in which an
intervention would be necessary in order to reach a well-founded group decision.

As discussed in Chapter~2, relevant knowledge in organizational collaboration is often distributed
asymmetrically across participants.
An adequate solution in a convergent decision process then depends on the group recognizing,
sharing, and integrating information that is initially held by only some members.
Hidden-profile experiments operationalize this asymmetry. There are different forms of hidden-profile experiments.
What these designs have in common is that no member's individual information set is sufficient to
identify the superior option; that option becomes apparent only when the distributed information is
disclosed and combined in the joint discussion.
In some designs, the information that is public from the outset is additionally constructed so as to
favor a suboptimal alternative; in others, this directional bias is not imposed.
In classic hidden-profile experiments in which the group has access to shared, given information
that suggests a inferior choice, while the split information suggests a better
choice, groups typically tend toward the inferior alternative. \cite{arkesInformationExchangeGroup2009}

The evaluation uses the hidden-profile hiring simulation published by Baker
\cite{bakerEnhancingGroupDecision2010}.
In this scenario, participants receive differently distributed information about three candidates,
and the information that is initially shared favors a suboptimal choice.
They must jointly decide whom they consider most suitable.
The task is therefore expected to elicit the interaction patterns a proactive moderator should
detect: premature consensus, unique information that remains unmentioned, unbalanced participation,
and a drift of the discussion towards commonly known talking points.
This specific scenario was selected for three further reasons:

\begin{itemize}
  \item \textbf{Feasibility in a remote meeting.}
    The task can be administered in groups of four and completed in a compact session of
    approximately 30 minutes.
    The materials can be distributed digitally, which allows the discussion to be conducted
    remotely.
  \item \textbf{Fit with the intended intervention situations.}
    The information structure of the task is expected to produce the patterns described above
    and thus occasions on which the prototype might intervene.
  \item \textbf{Replicability.}
    Baker provides the full experimental materials with the article, including role information,
    participant-specific information sheets, and evaluation forms, so that the scenario can be
    reused in the present study.
\end{itemize}

Meeting success was assessed with the questionnaire instrument revalidated by Davison for the
systematic measurement of meeting outcomes
\cite{davisonInstrumentMeasuringMeeting1999}.
The instrument covers five dimensions: communication, quality of discussion, status effects
(power structures and hierarchies), teamwork, and efficiency.
These dimensions correspond to the aspects of group interaction that the hidden-profile task is
intended to place under strain.

The technology-related items of the instrument, however, were formulated for Group Support
Systems of that period, that is, for largely passive, computer-supported collaboration tools.
They do not address autonomous, proactive AI agents.
In addition, Houtti et al.\ found that participants rated meetings with their virtual co-host
more positively overall, yet tended to rationalize or ignore its critical feedback
\cite{houttiObserveAskIntervene2025}.
Items on interaction with a proactive AI agent were therefore added:

\begin{itemize}
  \item \textbf{Intrusiveness and timing.}
    The AI agent's contributions came at moments when they did not disrupt the conversational
    flow but supported it.
  \item \textbf{Context-appropriateness.}
    The agent seemed to understand what we were currently discussing.
  \item \textbf{AI-mediated social comfort.}
    It was more comfortable for me to receive this feedback from the AI than if another group
    member had pointed it out to me directly.
\end{itemize}

\subsection{Sample}
\subsection{Data Collection}
